\section{Заключение}
\label{sec:conclusion}
SKILL.state предлагает полезное отделение рабочей памяти от растущего транскрипта. Перенос этой идеи в программирующего агента подтверждает возможность уменьшить пиковый контекст, но показывает, что сжатие само по себе не сохраняет успешность и не гарантирует снижения расходов.

Состав состояния меняет результат. В поисковом сравнении автоматические записи изменений и команд заметно лучше бедной памяти со ссылками на файлы; последующее обогащение не даёт монотонного улучшения. В повторных измерениях сохранение недавней истории восстанавливает значительную часть наблюдаемого разрыва и для состояния модели, и для гибридной памяти. Проверенные изменения схемы и протокола такого восстановления не обеспечивают, а экономический эффект истории зависит от реализации.

Практический вывод состоит в разделении решений о памяти: какие наблюдения сохраняет среда исполнения, какую интерпретацию поддерживает модель и какие исходные взаимодействия остаются доступными. Эти решения следует оценивать по успешности и стоимости всего выполнения наряду с размером контекста. Испытанные способы частично компенсируют ограничения замены истории состоянием; их результаты не устанавливают ни универсального преимущества гибридной архитектуры, ни дополнительной пользы состояния перед одним окном недавней истории.
